% Standalone section material; the parent article supplies theorem environments.
% Bibliography keys are supplied in su2_sources.bib.
% Source audit and suggested entries appear at the end as LaTeX comments.

\section{An exact algebraic route controlled by presentation size}
\label{sec:su2-presentation}

The reduction from knottedness to real algebraic feasibility is established
work: Meesum and Prathamesh gave a constant-degree formulation with four real
variables per Wirtinger generator \cite{MP2018}.  The issue addressed here is
how to reduce the number of independent quaternion variables without hiding
an exponential increase in polynomial degree or coefficient size.  We prove
three quantitative compilation results.  They concern explicit words,
straight-line word circuits, and a hybrid representation that retains selected
circuit nodes.  None presupposes a fast algorithm for finding an optimally
small presentation.

All complexity statements below refer to deterministic bit operations.
The notation $A^{O(B)}$ has a universal implicit constant.  A polynomial
factor in an input length always has a universal exponent independent of the
number of generators. A knot-group presentation is supplied with a verified
derivation, or an explicit knot-group promise when stating a conditional
theorem. The cost of establishing that origin is discussed separately
in Section~\ref{subsec:su2-provenance}.

\subsection{The topological and real-algebraic ingredients}

The only deep knot-theoretic input needed in this section is the following
consequence of Kronheimer--Mrowka.

\begin{theorem}[Kronheimer--Mrowka consequence]
\label{thm:su2-detection}
For a tame knot $K\subset S^3$,
\[
 K\text{ is nontrivial}
 \quad\Longleftrightarrow\quad
 \pi_1(S^3\setminus K)\text{ admits a homomorphism to }\mathrm{SU}(2)
 \text{ with nonabelian image}.
\]
\end{theorem}

\begin{proof}
The unknot group is infinite cyclic, so every image of it is abelian.
For the other direction, apply \cite[Theorem~1]{KM2004SU2} to $+1$ surgery
on a nontrivial knot.  The resulting fundamental group has a homomorphism
to $\mathrm{SU}(2)$ with noncyclic image.  Its first integral homology is
zero, so an abelian image would be trivial.  Consequently the image is
nonabelian.  Composing with the quotient map from the knot group gives the
required representation.  This argument explains why the distinction between
``noncyclic'' and ``nonabelian'' in the cited theorem causes no difficulty.
\end{proof}

We use the following established form of existential real decision.

\begin{theorem}[Existential real decision bound]
\label{thm:su2-real-decision}
Let a quantifier-free formula of bounded Boolean length involve a fixed number of integer polynomials
in $k\geq 1$ real variables, each of total degree at most $D\geq 2$, with
coefficient bit length at most $\tau$.  Existence of a real solution can be
decided deterministically in
\[
 \operatorname{poly}(\tau+k)\,D^{O(k)}
\]
bit operations, with the dense polynomial input included in the bound.
Strict inequalities are allowed.
\end{theorem}

This is a standard consequence of Renegar's existential decision theorem
\cite{Renegar1992}.  An independently accessible account is
\cite[Theorem~2.18]{BasuSurvey}: for one quantified block and no free
variables, the arithmetic-operation bound is $(sD)^{O(k)}$, and the bit
length of intermediate integers is at most $\tau D^{O(k)}$.  Here $s$ is
the number of input polynomials.  Using ordinary integer arithmetic and
fixing $s$ gives the displayed bit bound.  The dense-input convention in
\cite[Section~2.5]{BasuSurvey} is essential.  The proofs below include the
cost of producing that dense input; a short arithmetic circuit alone is not
treated as a short dense polynomial.

\subsection{Unit quaternions and the correct nonabelian test}

Write a quaternion as $q=(a,\mathbf v)$, with $a\in\mathbb R$ and
$\mathbf v\in\mathbb R^3$.  Its product and conjugate are
\begin{align}
 (a,\mathbf v)(b,\mathbf w)
   &=\bigl(ab-\mathbf v\cdot\mathbf w,
       a\mathbf w+b\mathbf v+\mathbf v\times\mathbf w\bigr),
       \label{eq:su2-product}\\
 \overline{(a,\mathbf v)}&=(a,-\mathbf v).
\end{align}
The unit quaternions, defined by $a^2+\|\mathbf v\|^2=1$, form
$\mathrm{SU}(2)$.  Conjugation gives inversion on this group.  In particular,
inverse letters require linear polynomials, not division.

Let $x_1,\ldots,x_r$ be the original generators of a presentation and let
$q_i=(a_i,\mathbf v_i)$ be their putative images.  Define
\begin{equation}
 H(q_1,\ldots,q_r)
   =\sum_{1\leq i<j\leq r}\|\mathbf v_i\times\mathbf v_j\|^2.
 \label{eq:su2-nonabelian}
\end{equation}
This is an integer polynomial of degree at most four.  The empty sum is zero.

\begin{lemma}[Nonabelian-image test]
\label{lem:su2-nonabelian}
A representation of a group generated by $x_1,\ldots,x_r$ has nonabelian
image if and only if $H>0$.
\end{lemma}

\begin{proof}
Equation~\eqref{eq:su2-product} gives
$q_iq_j-q_jq_i=(0,2\mathbf v_i\times\mathbf v_j)$.
Thus $q_i$ and $q_j$ commute exactly when their cross product vanishes.
A group generated by a pairwise commuting set is abelian.  Finally, a sum
of real squares is positive exactly when at least one of its summands is
nonzero.
\end{proof}

\begin{lemma}[Conjugate meridional generators]
\label{lem:su2-meridional}
Suppose the generators $x_1,\ldots,x_r$ are pairwise conjugate in the
presented group.  For every representation their scalar quaternion parts
are equal, and the lower-degree polynomial
\[
 H_{\mathrm{mer}}=\sum_{i=2}^r\|q_i-q_1\|^2
\]
is positive exactly when the image is nonabelian.
\end{lemma}

\begin{proof}
Quaternion conjugation preserves the scalar part, so the asserted scalar
equalities hold.  If the entire image is abelian, the images of the words
conjugating $x_1$ to the other generators commute with $q_1$; consequently
all generator images equal $q_1$.  Conversely, if all generator images
are equal, they generate a cyclic, hence abelian, subgroup.  The
sum-of-squares test expresses the negation of these equalities.
\end{proof}

The hypothesis concerns conjugacy \emph{in the presented group}, not just
the fact that selected quaternion values happen to have the same trace.
For a complete Wirtinger presentation of a knot, or a certified generating
set of consistently oriented meridians obtained from it, it applies.
That shortcut does not survive arbitrary changes of generators.  For example,
\[
 \langle a,b\mid b=a^2\rangle\cong\mathbb Z
\]
is an unknot-group presentation.  Sending $a$ to the quaternion $i$ and $b$
to $-1$ gives distinct generator images and still gives an abelian image.
The test~\eqref{eq:su2-nonabelian} remains correct for this presentation.
Auxiliary variables introduced during compilation are word values, not
additional independent generators; it suffices to test the original $r$
generators.

\begin{lemma}[A universal simultaneous-conjugation gauge]
\label{lem:su2-gauge}
Every tuple $q_1,\ldots,q_r$ of unit quaternions with $r\geq2$ is
simultaneously conjugate to a tuple of the form
\[
 q_1=(a_1,b_1,0,0),\qquad q_2=(a_2,b_2,c_2,0),
\]
with the later quaternions unrestricted.  The assertion includes central
$q_1$, central $q_2$, and parallel imaginary parts.  For $r=1$, the
first displayed condition alone is always attainable.  Hence imposing
these coordinate zeros loses no solution to the representation and
nonabelian-image conditions above.
\end{lemma}

\begin{proof}
Let $\mathbf v_1,\mathbf v_2$ be the first two imaginary parts.  If
$\mathbf v_1\ne0$, choose the first vector of a positively oriented
orthonormal frame to be $\mathbf v_1/\|\mathbf v_1\|$.
If the projection of $\mathbf v_2$ to its orthogonal plane is nonzero,
choose the second frame vector in that projection's direction; otherwise
choose any perpendicular unit vector.  The third vector is the cross
product of the first two.  Coordinates in this frame put $\mathbf v_1$
on the first axis and $\mathbf v_2$ in the first two coordinate directions.
If $\mathbf v_1=0$, choose the first frame vector along $\mathbf v_2$
when it is nonzero, and choose any positively oriented orthonormal frame
when both vectors vanish.  This also has the required properties.  The
case $r=1$ uses just the first-vector construction.

The resulting coordinate change is in $\mathrm{SO}(3)$.  Every such
rotation is induced by conjugation by a unit quaternion: rotation through
angle $\theta$ about the unit vector $\mathbf n$ is realized by
$u=(\cos(\theta/2),\mathbf n\sin(\theta/2))$, as follows by
substituting in~\eqref{eq:su2-product}.  Replace every generator and
every auxiliary word value by $uqu^{-1}$.  Conjugation preserves unit
norms, word equations, and commutation or equality of generator images.
The scalar parts are unchanged.  Thus all constraints are preserved.
\end{proof}

No coordinate is required to be strictly positive or nonzero in this
gauge.  It removes three real generator coordinates when $r\geq2$,
and two when $r=1$.  If the conjugate-meridian hypothesis also holds,
one may additionally share a single scalar variable.  For $r\geq2$
the generator tuple then uses $3r-2$ real variables; the auxiliary word
variables still use four each.  The general bounds below use the simpler
upper bound of four variables per quaternion, so these safe reductions
can only improve them.

\begin{theorem}[Known traceless-meridian detection]
\label{thm:su2-traceless}
A knot $K\subset S^3$ is nontrivial if and only if its group has a
nonabelian $\mathrm{SU}(2)$ representation carrying a chosen meridian
to the quaternion $i=(0,1,0,0)$.
\end{theorem}

The forward implication is the established theorem of Kronheimer--Mrowka,
\cite[Corollary~7.17]{KM2010Excision}; the reverse implication follows
because the unknot group is cyclic. ``Irreducible'' in that source is
equivalent to nonabelian image for representations into $\mathrm{SU}(2)$.
This is a stronger published ingredient than the unrestricted detection
consequence used to prove the general presentation bounds.

\begin{corollary}[A smaller formula for certified meridians]
\label{cor:su2-fixed-meridian}
For a knot-group presentation on $r\ge2$ certified consistently oriented
meridians, exact nontriviality can be tested while imposing
\[
 q_1=(0,1,0,0),\qquad q_2=(0,b_2,c_2,0),\qquad
 q_j=(0,b_j,c_j,d_j)\quad(j\ge3).
\]
With $u$ retained auxiliary quaternions this uses $3r-4+4u$ real
variables. The degree and coefficient estimates of the subsequent
compilation theorems do not increase.
\end{corollary}

\begin{proof}
Theorem~\ref{thm:su2-traceless} fixes the first meridian image to $i$.
Every other meridian is conjugate to it in the knot group, so all scalar
parts vanish. A remaining rotation about the first imaginary axis puts
the second imaginary vector in the first two coordinate directions,
including the parallel case. The second quaternion contributes two
variables, the later $r-2$ each contribute three, and each auxiliary
contributes four. Substitution of zero and one cannot increase total
degree or coefficient $\ell^1$ norm, so the coefficient-height estimates
proved below remain valid.
\end{proof}

This corollary is not a new trace-existence theorem. It describes a
sound further specialization of the meridional compiler. The recorded
implementation retains its general gauge and uses trace zero only as a
SAT-witness probe; it does not implement this smaller complete formula.
Consequently the reported dimensions and timings below are not those of
the stronger specialization. For arbitrary generators, a chosen generator
need not represent a meridian, and neither the trace restriction nor the
fixed value $i$ is justified without that additional provenance.

\subsection{Explicit relators: a degree--variable tradeoff}

Let
\[
 \mathcal P=\langle x_1,\ldots,x_r\mid w_1,\ldots,w_m\rangle,
 \qquad r\geq 1,
\]
be a finite presentation.  The relators are explicit words in the generators
and their inverses, and
\[
 N=\sum_{j=1}^m |w_j|
\]
is their total letter length.  Empty relators are discarded.  This standard
normalization ensures $m\leq N$ if $N>0$; otherwise there are no relators.
Encoding a generator index costs $O(\log(r+1))$ bits, so the input length
is polynomial in $N+r$.  If an input format permits an arbitrarily long
list of empty relators, its literal parsing time must be added separately.

\begin{theorem}[Explicit-presentation tradeoff]
\label{thm:su2-explicit-tradeoff}
For any integer $L\geq 1$, existence of a nonabelian $\mathrm{SU}(2)$
representation of $\mathcal P$ can be decided in
\begin{equation}
 \operatorname{poly}(N+r)\,(L+2)^{O(r+N/L)}
 \label{eq:su2-explicit-tradeoff}
\end{equation}
bit operations, where one may restrict $L\leq\max(1,N)$.
In particular, it can be decided in
\begin{equation}
 \operatorname{poly}(N+r)
 \left(2+\frac{N}{r}\right)^{O(r)}.
 \label{eq:su2-explicit-optimized}
\end{equation}
If $\mathcal P$ presents the group of a supplied knot $K$, these are
also bounds for recognizing whether $K$ is the unknot, with the
presentation-provenance cost added.
\end{theorem}

\begin{proof}
Partition each nonempty relator $w_j$ into consecutive chunks
\[
 w_j=c_{j,1}\cdots c_{j,h_j},\qquad
 h_j=\lceil |w_j|/L\rceil,\qquad 1\leq |c_{j,a}|\leq L.
\]
Introduce a quaternion variable $y_{j,a}$ for each intermediate prefix
$c_{j,1}\cdots c_{j,a}$, with $1\leq a<h_j$.  No such variable is needed
for the empty prefix or the final prefix, whose values are prescribed to be
$1$.  The number $u$ of auxiliary quaternion variables satisfies
\[
 u=\sum_j(h_j-1)\leq N/L.
\]
There are consequently $k=4(r+u)$ real variables.

Let $C_{j,a}$ denote the quaternion polynomial obtained by evaluating the
letters of $c_{j,a}$ using the original variables $q_i$ and their
conjugates.  Each coordinate of $C_{j,a}$ has degree at most $L$.
Impose the quaternion equations
\[
 y_{j,1}=C_{j,1},\qquad
 y_{j,a}=y_{j,a-1}C_{j,a}\quad(1<a<h_j),\qquad
 1=y_{j,h_j-1}C_{j,h_j}.
\]
When $h_j=1$, use just $1=C_{j,1}$.  Each quaternion equation means four
real equations, of degree at most $L+1$.  Also impose the unit-norm equation
on every original and auxiliary quaternion.  The latter equations are
redundant for correctly evaluated auxiliary variables, but their inclusion
is harmless and makes the intended domain explicit.

For any assignment of unit quaternions to the original generators, the
prefix equations have a unique assignment to all the auxiliary variables.
This assignment satisfies the final equation for $w_j$ if and only if
$w_j(q_1,\ldots,q_r)=1$.  This proves that the equations encode precisely
the representations of the original presented group.

List the scalar polynomial residuals of all equations as
$p_1,\ldots,p_A$ and set
\[
 F=\sum_{\nu=1}^{A}p_\nu^2.
\]
Over the real numbers, $F=0$ is equivalent to all the original equalities.
The required feasibility formula therefore has only two polynomial atoms:
\begin{equation}
 \exists(q,y)\quad F(q,y)=0\ \wedge\ H(q)>0.
 \label{eq:su2-two-atoms}
\end{equation}
Its maximum degree is at most $D=2L+2$, which is at least four.

We now account for bit size and dense compilation.  A coordinate of a
product of $\ell$ quaternion letters has coefficient $\ell^1$ norm at
most $4^{\ell-1}$ for $\ell\geq1$: at each multiplication, one coordinate
is a signed sum of four products.  Conjugation only changes signs.
An auxiliary prefix factor increases the word length by at most one.
Thus each residual has coefficient $\ell^1$ norm $2^{O(L)}$, and squaring
it preserves this form of bound.  There are $O(N+r)$ residuals.  Every
coefficient of $F$ or $H$ consequently has bit length
\[
 \tau=O\bigl(L+\log(N+r+2)\bigr).
\]
This argument allows coefficients to accumulate when identical monomials
are collected; it does not assume that aggregate coefficients remain in a
fixed finite set.

A polynomial of degree at most $D$ in $k$ variables has at most
\[
 \binom{k+D}{k}\leq(D+1)^k
\]
monomials.  Expand each chunk by successive multiplication by linear
quaternion letters, multiply by its optional prefix variable, square the
residuals, and collect equal monomials using exact integer arithmetic.
Even naive dense multiplication costs at most a polynomial in the number
of monomials.  All dense polynomials can therefore be constructed in
\[
 \operatorname{poly}(N+r+L)\,(D+1)^{O(k)}
\]
bit operations.  For $L\leq\max(1,N)$ this has the form
\eqref{eq:su2-explicit-tradeoff}.  Applying
Theorem~\ref{thm:su2-real-decision} to~\eqref{eq:su2-two-atoms} has the
same bound.  This includes both compilation and decision.

Finally take $L=\max(1,\lceil N/r\rceil)$.  Then $u\leq r$ when
$N\geq r$, while $u\leq N\leq r$ when $N<r$, and
$L+2=O(2+N/r)$.  This gives~\eqref{eq:su2-explicit-optimized}.
When $N=0$ there are only the original quaternion variables and the same
argument with $L=1$ applies.  The knot corollary follows from
Theorem~\ref{thm:su2-detection}.
\end{proof}

\begin{remark}[What the tradeoff means]
The two extreme formulations are $L=1$, with constant degree and
$O(N+r)$ quaternion variables, and $L=N$, with no prefix variables and
degree $O(N)$.  Introducing some prefix variables can improve the
asymptotic bound over complete substitution.  The optimized choice is a
guarantee, not a recommendation to use that value in every implementation:
the actual support of the polynomials and the individual relator lengths
may justify a different value.
\end{remark}

\subsection{A concrete consequence for the supplied diagram}

For a nonempty oriented knot diagram $D$, traverse its unique component
and record an $O$ at an overcrossing encounter and a $U$ at an
undercrossing encounter.  Define $b(D)$ to be the number of maximal
$O$-runs in this cyclic word.  This is the number of maximal diagrammatic
overpasses, not the minimum bridge number of the knot type.

\begin{lemma}[Overpass presentation]
\label{lem:su2-overpass}
From an $n$-crossing knot diagram with $b=b(D)$, one can construct a
presentation of its knot group with $b$ meridional generators and total
explicit relator letter length at most $2n+2b$.  The construction and a
check of the elimination steps take polynomial time in $n$.
\end{lemma}

\begin{proof}
Choose a meridian $x_i$ for each maximal overpass.  Along the intervening
underpass, each crossing conjugates the current meridian by the meridian
of the overpassing strand, with the sign determined by the oriented
crossing convention.  If the successive conjugating letters are
$z_1,\ldots,z_\ell$, all of them are among the selected overpass
generators and their inverses.  Eliminating the intermediate Wirtinger
generators gives one relation
\[
 x_{\mathrm{next}}=w x_{\mathrm{previous}}w^{-1},\qquad |w|=\ell,
\]
with the appropriate order of factors.  This elimination uses only
relations that define a new intermediate generator from preceding data,
so it is a sequence of valid Tietze eliminations.  The underpasses partition
the $n$ undercrossing encounters, hence $\sum\ell=n$.  Writing each of
the $b$ resulting relations as a relator gives total length
$\sum(2\ell+2)=2n+2b$.
\end{proof}

\begin{corollary}[Diagrammatic-overpass bound]
\label{cor:su2-overpass}
Unknot recognition on a supplied $n$-crossing diagram $D$ can be performed
deterministically in
\[
 \operatorname{poly}(n)\left(2+\frac{n}{b(D)}\right)^{O(b(D))}.
\]
In particular, the bound is polynomial for bounded $b(D)$ and is
$n^{O(\log n)}$ when $b(D)=O(\log n)$.  For arbitrary diagrams it is
at most $2^{O(n)}$ times a polynomial factor.
\end{corollary}

\begin{proof}
Apply Theorem~\ref{thm:su2-explicit-tradeoff} and
Lemma~\ref{lem:su2-overpass}.  The last assertion follows from
$b\log(2+n/b)=O(n)$ for $1\leq b\leq n$.
\end{proof}

This result is deliberately parameterized by the supplied diagram.
An alternating diagram has cyclic encounter word $OUOU\cdots OU$, so
$b(D)=n$, even when its knot type has bridge number two.  The usual
alternating diagram of $T(2,m)$ is an example.  The corollary does not
provide a polynomial-time conversion to a bridge-minimizing diagram.
It also does not replace $b(D)$ by the number of strands of a supplied
braid.  Iterated Artin substitutions can produce very long expanded words;
the next theorem is the appropriate bound when these substitutions are
kept as circuits.

\begin{remark}[The optimized bridge compiler]
The implementation need not write both copies of every conjugating word.
It may instead impose $q_{\mathrm{next}}q_w=q_wq_{\mathrm{previous}}$.
Unit generator norms and the prefix definitions ensure that $q_w$ is
invertible, so this is precisely the same relation.  Split the $n$
conjugator letters into chunks of at most $L$.  There are
$u\leq n/L$ prefix quaternion variables.  The final conjugacy equation
has degree at most $L+2$: one prefix factor, at most $L$ final letters,
and one endpoint generator.  Thus the aggregate degree is at most
$2L+4$, while coefficient bit lengths remain
$O(L+\log(n+b+2))$.  Lemmas~\ref{lem:su2-meridional} and
\ref{lem:su2-gauge} reduce the number of real variables to
$3b-2+4u$ for $b\geq2$.  When $b=1$, the nonabelian test is identically
zero and the decision is immediate.  This more economical compiler
retains the asymptotic bound of Corollary~\ref{cor:su2-overpass}.
\end{remark}

\subsection{Compressed words: retaining every circuit gate}

A straight-line word circuit is a finite directed acyclic graph whose
inputs are $x_1,\ldots,x_r$ and, optionally, the identity, and whose
internal gates are concatenation and inversion.  A list of designated
outputs specifies the relators.  The circuit may share subwords.  Let $s$
be the number of internal gates after expanding any binary integer-power
instruction into repeated squaring and multiplication.  Negative powers
also use inversion.  Thus a power $w^e$ costs $O(\log(|e|+1))$ gates;
it is not counted as one constant-degree gate regardless of $e$.
Let $S$ be the bit length of the whole circuit input, including all output
references and exponent encodings.

\begin{theorem}[Straight-line presentation bound]
\label{thm:su2-slp}
For a presentation specified by a straight-line word circuit with $r$
generators and $s$ elementary internal gates, existence of a nonabelian
$\mathrm{SU}(2)$ representation can be decided in
\[
 \operatorname{poly}(S+r+s)\,2^{O(r+s)}
\]
bit operations.  The feasibility formula has $4(r+s)$ real variables,
one equality and one strict inequality, and maximum degree four.
Its aggregate coefficient bit length is $O(\log(S+r+s+2))$.
\end{theorem}

\begin{proof}
Assign a quaternion variable to every generator and every internal gate.
For a concatenation gate $z=uv$, impose the four coordinate equations
$q_z-q_uq_v=0$.  They have degree at most two.  For an inversion gate
$z=u^{-1}$, impose $q_z-\overline{q_u}=0$, which is linear.
Identity inputs are constants.  Impose $q_z=1$ for each relator output
and unit-norm equations for all quaternion variables.  Induction in a
topological ordering of the circuit proves that every gate variable has
the unique correct value for the chosen generator images.  The equations
therefore encode exactly the original group representations, regardless
of how long the expanded words would be.

Square and sum all scalar residuals to form $F$, and retain
$H>0$ from~\eqref{eq:su2-nonabelian}.  Both $F$ and $H$ have degree at
most four.  Every residual has a bounded number of monomials and bounded
integer coefficients, even when a gate repeats the same input.  Collecting
terms from $O(S+r+s)$ residuals, and from $O(r^2)$ terms in $H$, gives
coefficient bit length $O(\log(S+r+s+2))$.
For fixed degree four, the number of monomials in $4(r+s)$ variables is
polynomial in $r+s$, so even full dense compilation is polynomial in the
input parameters.  Theorem~\ref{thm:su2-real-decision} then gives the bound.
\end{proof}

\begin{remark}[A circuit gate is not a free degree reduction]
The chain $z_0=x$, $z_{j+1}=z_jz_j$ computes $x^{2^s}$ with $s$ gates.
Eliminating every gate can leave coordinate polynomials of formal degree
$2^s$.  Retaining the gates gives quadratic local constraints but adds
four real variables per gate.  Theorem~\ref{thm:su2-slp} pays for those
variables.  It does not infer quasipolynomial time from a small number of
original generators when the circuit has linear size in the crossing count.
\end{remark}

\subsection{A hybrid bound for a cut in the word circuit}

Full expansion and full gate retention are special cases of one construction.
Choose $t$ internal gates to retain as quaternion variables.  For every
retained gate, expand its defining word through unretained gates until
another retained gate or an original generator is reached.  Do the same
for every relator output.  In this expansion, inversion reverses the word
and inverts its letters.  Identity letters are discarded.  Let $d\geq1$
bound the length of every resulting word in the cut variables and their
inverses.  Equivalently, it is enough to bound the formal degree obtained
by assigning degree one to generator and cut variables, degree zero to
identity, taking sums at concatenation gates, and preserving degree at
inversion gates.  For a retained gate's own definition, apply its operation
before stopping at upstream cut gates.

\begin{theorem}[Hybrid cut bound]
\label{thm:su2-hybrid}
With the notation above, a supplied choice of $t$ retained gates and degree
bound $d$ gives a deterministic decision algorithm with bit complexity
\[
 \operatorname{poly}(S+r+s)\,(d+2)^{O(r+t)}.
\]
The cost of expanding and collecting all relevant coordinate polynomials
is included.  The theorem is for word circuits with concatenation,
inversion, and identity constants; it does not assert the same
coefficient-height bound for unrestricted arithmetic circuits.
\end{theorem}

\begin{proof}
Use quaternion variables only for the original generators and the retained
gates.  Each retained gate is constrained to equal the polynomial word in
upstream cut variables given by its defining expansion.  Each relator word
is constrained to be the identity.  Add unit-norm equations and the
nonabelian test on the original generators.  Topological induction over
retained gates proves equivalence with the full circuit equations: there
are no circular definitions, and unretained gates are uniquely evaluated
from their inputs.

After summing squares, there are two polynomial atoms in
$k=4(r+t)$ variables, with maximum degree at most
$\max(4,2d)$.  Every expanded expression is a genuine word of length at
most $d$ in quaternion variables or conjugates.  Its coefficient
$\ell^1$ norm is at most $4^{d-1}$, so the aggregate coefficient bit
length is $O(d+\log(S+r+s+2))$.  In particular, repeated squaring of
arbitrary integer constants cannot create an unaccounted height explosion:
such constants and operations are not part of the word-circuit language.

Memoize the polynomial of each relevant unretained gate, treating each
retained gate as a new variable when used by a later gate.  Store its
defining polynomial separately.  All intermediate nonconstant formal
degrees are at most $d$, since concatenation adds nonnegative degrees
and inversion preserves them.  A dense polynomial of degree $2d+4$ in
$k$ variables has at most $(2d+5)^k$ monomials.  Ordinary exact dense
multiplication, summation, and coefficient collection therefore cost
\[
 \operatorname{poly}(S+r+s+d)\,(d+2)^{O(k)}.
\]
The same bound covers existential real decision.  Since $r\geq1$, every
fixed polynomial factor in $d$ is absorbed into $(d+2)^{O(r+t)}$.
This gives the claimed form, including dense compilation.
\end{proof}

A useful, fully computable size objective for a proposed cut is
\[
 \Lambda=(r+t)\log_2(d+2).
\]
It reflects the exponent in the proved worst-case bound, while practical
decisions can additionally use monomial counts and coefficient heights.
Optimizing $\Lambda$ over all cuts is a separate combinatorial problem;
the theorem assumes the cut is supplied.  One can always evaluate a fixed
list of cuts and keep the best proved estimate, or run their solvers in a
fair portfolio.  No exhaustive search over cuts is required for correctness.

\begin{remark}[A greedy threshold gives a certified cut]
A simple implementable rule processes a concatenation DAG in topological
order.  Fix a threshold $h\geq1$.  Store degree at most $h$ for every
currently expanded value.  At a product of two stored values, compute the
sum of their formal degrees.  If it exceeds $h$, retain the product gate,
record its defining equation, and reset its stored degree to one;
otherwise leave it expanded.  Every recorded gate definition has degree
at most $2h$, and every output value has degree at most $h$.  Consequently
the rule supplies a valid cut with $d\leq2h$ in
Theorem~\ref{thm:su2-hybrid}, and its aggregate formula degree is at
most $4h$.  The threshold is a bound on stored expressions, not on a
retained gate's defining polynomial.  For $h=1$, all nontrivial products
are retained and the aggregate degree is at most four.  This argument
also accounts for the temporary multiplication performed before deciding
to retain a gate; that multiplication never exceeds degree $2h$.
\end{remark}

\begin{corollary}[A sufficient quasipolynomial compilation condition]
\label{cor:su2-qp-condition}
Suppose an $n$-crossing knot diagram can be converted, with verified
group-preserving provenance, to a word circuit and a cut in
$\operatorname{poly}(n)$ time.  Suppose the circuit description has
polynomial length and
\[
 (r+t)\log(d+2)=O((\log n)^2).
\]
Then this conversion followed by exact existential real decision recognizes
the unknot in $n^{O(\log n)}$ time.  In particular, it is enough to have
$r+s=O((\log n)^2)$ with every gate retained.
\end{corollary}

This is a sufficient condition on an explicitly constructed representation.
No assertion that every knot diagram admits such a construction is proved
here.  In fact, requiring a small presentation of the \emph{entire} knot
group is too restrictive for a universal route, as the following elementary
obstruction shows.

\begin{proposition}[A rank obstruction to universal presentation compression]
\label{prop:su2-rank-obstruction}
For every $m\geq1$, there is a knot with a $3m$-crossing diagram whose
group has minimum generator number exactly $m+1$.  Consequently no
group-preserving compiler can satisfy
$(r+t)\log(d+2)=O((\log n)^2)$ on every $n$-crossing diagram solely
by choosing a presentation of its entire knot group and a cut in its
relator circuit.
\end{proposition}

\begin{proof}
Let $K_m$ be the connected sum of $m$ trefoils, drawn by joining $m$
three-crossing trefoil diagrams without introducing new crossings.
The Wirtinger presentation of a trefoil is
$\langle a,b\mid aba=bab\rangle$, with $a,b$ meridians.
Identifying the meridian $a$ in the connected sum, either by the
Wirtinger construction or van Kampen's theorem, gives
\[
 G_m=\langle a,b_1,\ldots,b_m\mid
     ab_i a=b_i a b_i\ (1\leq i\leq m)\rangle.
\]
Let $V=\mathbb F_3^m$ with standard basis $e_1,\ldots,e_m$, and
let $\operatorname{Dih}(V)=V\rtimes C_2$ with multiplication
\[
 (v,\epsilon)(w,\delta)
  =\bigl(v+(-1)^\epsilon w,\epsilon+\delta\bmod2\bigr).
\]
The assignment $a\mapsto(0,1)$ and $b_i\mapsto(e_i,1)$ respects
the relators: the images of $ab_i a$ and $b_i a b_i$ are respectively
$(-e_i,1)$ and $(2e_i,1)$, and these coincide over $\mathbb F_3$.
It is surjective because $ab_i$ maps to $(-e_i,0)$, and these elements
generate $V$.

The group $\operatorname{Dih}(V)$ requires $m+1$ generators.  To
see the lower bound, start with any generating set of size $r$ and
choose one generator $z$ outside $V$.  Such an element exists and has
order two.  Replace each other generator outside $V$ by its product
with $z$; retain the generators already in $V$.  These replacements
preserve the generated subgroup and leave $z$ together with at most
$r-1$ vectors in $V$.  Conjugation by $z$ negates every vector of
$V$.  Therefore the rotation subgroup generated by this set is just
the $\mathbb F_3$ span of those at most $r-1$ vectors.  To obtain
all of $V$ requires $r-1\geq m$.  The displayed quotient then gives
$d(G_m)\geq m+1$, and the original presentation gives the reverse
inequality.  Since $t\geq0$ and $d\geq1$,
the proposed cut objective is then $\Omega(m)=\Omega(n)$ on this
family, regardless of relator compression.
\end{proof}

This obstruction concerns the requirement to preserve the whole group,
not the complexity of deciding the unknot.  Connected sums admit a
decision decomposition: a connected sum is trivial exactly when all its
summands are trivial.  A universal strategy using the algebraic route
would therefore need to exploit such certified decompositions, or other
logically justified reductions, before demanding small presentation
parameters.  Finding and certifying a decomposition also has a cost.

\begin{proposition}[Prime knots also obstruct universal rank compression]
\label{prop:su2-prime-rank-obstruction}
For every odd $m\geq3$, the pretzel knot
$P_m=P(3,\ldots,3)$ with $m$ entries has a prime alternating
$3m$-crossing diagram, and its knot group requires at least $m$
generators.
\end{proposition}

\begin{proof}
Use the standard pretzel diagram consisting of a cyclic row of $m$
vertical three-half-twist boxes, with adjacent upper endpoints joined
above the boxes and adjacent lower endpoints joined below them.  Each
box exchanges its two strands.  Traversing through two successive boxes
advances the upper connecting-arc index by two modulo $m$.  Since $m$
is odd, this visits every upper connecting arc and the diagram has one
component.  It has $3m$ crossings and is alternating.  One of its
checkerboard graphs consists of $m$ internally disjoint paths of length
three between the same two vertices.  This graph has no loop and no cut
vertex.  Hence its medial diagram is prime: a circle meeting the medial
diagram twice and having crossings on both sides would give either a
loop or a separating vertex of the checkerboard graph.  The theorem
that a prime alternating diagram represents a prime link, due to
Menasco~\cite{Menasco1984}, now gives primeness; see also the extension proved in
\cite[Theorem~10.3]{PurcellTsvietkova2023}, whose hypotheses specialize
to a prime alternating diagram on $S^2\subset S^3$.

We construct a quotient of the knot group.  Fox colorings with values
in a vector space $V$ over $\mathbb F_3$ satisfy
$v_{\mathrm{out}}=2v_{\mathrm{over}}-v_{\mathrm{in}}$ at each
undercrossing.  With one consistent choice of two-strand crossing
convention, the transformation of a pair of colors through a half twist
is
\[
 T=\begin{pmatrix}2&-1\\1&0\end{pmatrix}=I+N,
 \qquad N^2=0.
\]
Thus $T^3=I$ over $\mathbb F_3$; reversing the half-twist convention
replaces $T$ by $T^{-1}$ and has the same consequence.  Every twist
box therefore preserves its pair of upper colors at the corresponding
lower positions.  The exterior upper and lower connections have the
same cyclic pattern.  Consequently the $m$ upper connecting arcs can
receive arbitrary values, and these values extend to all arcs of the
diagram without further constraints.

Take $V=\mathbb F_3^{m-1}$ and assign the $m$ upper connecting arcs
the values $0,e_1,\ldots,e_{m-1}$.  Map each Wirtinger meridian
with color $v$ to the reflection $(v,1)$ of
$\operatorname{Dih}(V)$.  At a crossing,
\[
 (w,1)(v,1)(w,1)=(2w-v,1),
\]
so every Wirtinger relation holds; the sign of the conjugating meridian
does not matter since its image is an involution.  The image contains
$(0,1)$ and $(e_i,1)$ for every $i$, hence contains all rotations and
is surjective.  The generator-rank argument in the preceding proof,
with $\dim V=m-1$, implies that this quotient, and therefore the knot
group, requires at least $m$ generators.
\end{proof}

Neither proposition is a lower bound on unknot-recognition time.  In
particular, these families have nonconstant Fox colorings that already
certify knottedness by elementary linear algebra.  The prime example
shows that prime decomposition alone cannot make whole-group generator
rank uniformly small.  The useful research target is instead a
portfolio of certified decision reductions, followed by controlled
presentation or circuit-cut parameters on the instances that remain.
The real-algebraic implication for any such residual instance is proved
above; establishing a universal bound for a residual class requires a
separate theorem about that class.
The rank arguments are included as explicit obstructions for this
compilation strategy; no priority claim is made for the knot-group rank
bounds or the use of dihedral quotients.  Lustig--Moriah's more general
ordinary-rank calculations for generalized Montesinos knots
\cite{LustigMoriah1993} are recalled in the author's account
\cite[Theorem~3.6]{Moriah2007}.  The elementary proof above makes the
specific obstruction needed here independent of that more general theory.

\subsection{Cyclic images, gauge choices, and exact solver outcomes}

Several implementation restrictions follow directly from the proofs.

First, the unit equations are indispensable for interpreting quaternion
conjugation as inversion.  In contrast, an additional norm equation at a
gate is logically redundant if all its ancestors are correctly constrained
and the original generators are unit quaternions.  Omitting those redundant
equations is safe only with that inductive invariant preserved.

Second, gauge fixing means simultaneous conjugation.  It can rotate
imaginary quaternion axes, but it preserves each generator's scalar part
and hence its trace. Setting a meridian equal to the identity is not a
gauge choice and destroys all nonabelian representations of a Wirtinger
knot group. Setting its scalar part to zero is also not a gauge choice,
but Theorem~\ref{thm:su2-traceless} makes it a complete restriction for
verified knot meridians. Fixing that meridian to $i$ is then legitimate
as in Corollary~\ref{cor:su2-fixed-meridian}. The present generic
solver interface uses only a trace-zero witness probe: a satisfying
model is decisive, and failure returns to the unrestricted query.
Common scalar parts may be imposed on a
certified consistently oriented meridional generating set, since those
generators are conjugate, but not on arbitrary generators or prefix gates.

Third, the decision formula is existential over the reals.  A satisfying
assignment certifies a nonabelian representation and therefore knottedness.
Infeasibility certifies the unknot only for a correctly constructed formula
with verified knot-group provenance.  A solver timeout, resource limit,
unknown result, or failure to find a numerical approximation is not an
infeasibility result.  Approximate solutions alone are not exact witnesses.

Finally, Theorems~\ref{thm:su2-explicit-tradeoff},
\ref{thm:su2-slp}, and~\ref{thm:su2-hybrid} use a theoretical
singly exponential existential real decision procedure.  A practical
implementation may instead pass the same exact integer-polynomial formula
to NLSAT, cylindrical algebraic decomposition, or another sound exact
backend.  This transfers the logical reduction, but does not transfer
Renegar's worst-case running-time bound to that implementation.  A benchmark
of the practical backend and the asymptotic existence theorem answer
different questions and must be reported separately.

\subsection{Presentation provenance and novelty boundaries}
\label{subsec:su2-provenance}

For a diagrammatic-overpass presentation, provenance is the explicit
Wirtinger elimination in Lemma~\ref{lem:su2-overpass}.  A more ambitious
presentation simplifier must supply comparably checkable evidence that its
output presents the same knot group.  A sequence of Tietze transformations
is one possible form.  Each retained or eliminated generator must have its
defining relation accounted for, and every change to a relator must be
justified.  The verifier must retain the correct relators and the correct
word multiplication order.  A homology calculation, equality of a few knot
invariants, or successful tests on a finite knot table does not establish
this group isomorphism.

If provenance costs $T_{\mathrm{prov}}$ bit operations to construct or
verify, add $T_{\mathrm{prov}}$ to the applicable bound.  A short
presentation that exists without a bounded-cost construction does not
give an algorithm with the same bound.  Likewise, a small number of
generators does not help if explicitly writing its relators has already
cost exponential time and no smaller circuit is supplied.

The underlying $\mathrm{SU}(2)$ reduction, the sum-of-squares aggregation,
and singly exponential real decision are established ingredients.  The
2018 paper explicitly suggests substitution to reduce variable count as a
possible direction for further work.  The present contribution is the
fully costed presentation-sensitive tradeoff, its concrete consequence for
the overpasses of the supplied diagram, and the circuit/cut formulations
that preserve word compression without concealing degree.  A targeted
literature search did not locate the exact bound
$\operatorname{poly}(N+r)(2+N/r)^{O(r)}$ for knot presentations, but
absence from that search is not a priority proof.  These results should be
presented as proved refinements of the established reduction, not as a
claim that the general unknot-recognition problem has been solved in a
new complexity class.

\subsection{Source corrections relevant to implementation}

The arXiv v1 text of \cite{MP2018} contains several local errors that do
not invalidate its intended reduction but should not be copied into code.
Its preliminary definition calls the matrices in $\mathrm{SU}(2)$
Hermitian; the intended word is unitary.  Its displayed quaternion-shaped
matrix parameterization and unit-norm condition use the correct group.
The declared convention for Algorithm~1 is that ``Yes'' means unknot,
but its satisfiable branch returns ``Yes'' for a formula that detects
knottedness.  Under that declared convention, the two output labels must
be interchanged.  Finally, after all squared residuals are added and like
terms are collected, integer coefficients need not stay in the fixed
small set claimed in its complexity discussion.  A logarithmic
coefficient-height bound in the number of residuals suffices for its
$2^{O(n)}$ conclusion and is the bound used above.

% Suggested bibliography entries / checked primary-source locations:
%
% KM2004SU2:
% P. B. Kronheimer and T. S. Mrowka,
% Dehn surgery, the fundamental group and SU(2),
% Mathematical Research Letters 11 (2004), 741--754.
% https://arxiv.org/abs/math/0312322
% https://arxiv.org/pdf/math/0312322
% Theorem 1, first page; choose slope +1 for the argument above.
%
% MP2018:
% Syed M. Meesum and T. V. H. Prathamesh,
% Unknot Recognition Through Quantifier Elimination,
% arXiv:1803.00413v1 (2018).
% https://arxiv.org/html/1803.00413v1
% Prop. 7 states the Renegar bit bound; Section 6 suggests substitutions.
% Hermitian/unitary typo: Section 2 first paragraph.
% Label reversal: Algorithm 1, declared result vs satisfiable branch.
% Aggregated coefficient-height error: Section 5.
%
% Renegar1992:
% James Renegar, On the computational complexity and geometry of the
% first-order theory of the reals. Part I: Introduction. Preliminaries.
% The geometry of semi-algebraic sets. The decision problem for the
% existential theory of the reals,
% Journal of Symbolic Computation 13(3) (1992), 255--299.
% https://doi.org/10.1016/S0747-7171(10)80003-3
% The bibliographic record was checked; full publisher page did not fetch.
%
% BasuSurvey:
% Saugata Basu, Algorithms in Real Algebraic Geometry: A Survey.
% Author-hosted primary account of the BPR algorithms:
% https://www.math.purdue.edu/~sbasu/raag_survey2011_final.pdf
% Theorem 2.18 (printed p. 13) gives both operation count and intermediate
% integer bit bounds; Section 2.5 (printed p. 14) states dense-input convention.
% This particular hosted version uses theorem number 2.18; newer versions
% may renumber. Prefer an explicit URL/version in the bibliography.
